% ======================= ÖZGEÇMİŞ =======================
\clearpage
\chapter*{ÖZGEÇMİŞ}
\addcontentsline{toc}{chapter}{ÖZGEÇMİŞ}

\noindent
\textbf{Kişisel Bilgiler}\\[4pt]
\begin{tabularx}{\textwidth}{@{}p{4cm}X@{}}
Adı Soyadı      & : Taha Berk TEREKLİ \\
Bölümü          & : Matematik Bölümü \\
Üniversite      & : Yıldız Teknik Üniversitesi, Fen-Edebiyat Fakültesi \\
\end{tabularx}

\vspace{0.6cm}
\noindent
\textbf{Eğitim}\\[4pt]
\begin{tabularx}{\textwidth}{@{}p{4cm}X@{}}
Lisans & : Yıldız Teknik Üniversitesi, Matematik Bölümü (2022--2026) \\
Sertifika & : Stanford University, Code in Place -- Computer Science (2023) \\
\end{tabularx}

\vspace{0.6cm}
\noindent
\textbf{Kısa Özgeçmiş}\\[4pt]
Taha Berk TEREKLİ, Yıldız Teknik Üniversitesi Matematik Bölümü lisans
öğrencisidir. Matematik eğitimi boyunca analitik düşünme, soyut cebirsel yapılar
ve bilgisayar destekli hesaplama konularına ilgi duymuş; bu ilgisini Python,
SQL ve veri analitiği alanlarındaki uygulamalı çalışmalarla desteklemiştir.

Profesyonel deneyimlerinde veri analizi, otomasyon ve iş süreçlerinin yazılımla
iyileştirilmesi üzerine çalışmıştır. Getir'de veri stajyeri olarak SQL ve Python
tabanlı veri çıkarımı, analiz ve raporlama süreçlerinde; Papara ve Latro'da ise
RPA ve otomasyon projelerinde görev almıştır. Ayrıca TÜBİTAK 2209-A kapsamında
derin öğrenme tabanlı tıbbi görüntü işleme üzerine bir araştırma projesinde yer
almaktadır.

\vspace{0.6cm}
\noindent
\textbf{Bitirme Tezi ile İlişkisi}\\[4pt]
Bu bitirme çalışması, yazarın matematiksel altyapısı ile Python ve sembolik
hesaplama deneyimini bir araya getirmektedir. Tez kapsamında kuantum gruplar,
temsil teorisi ve Yang--Baxter tipi bağıntılar Python/SymPy ortamında
modellenmiş; geliştirilen \texttt{quantum\_group} paketi ve test takımı açık
kaynak olarak \url{https://github.com/TerekliTahaBerk/quantum-groups}
adresinde yayımlanmıştır.

\end{document}
